% !TeX root = ../main.tex
\section{Architecture Decisions}\label{sec:architecture}

In this section, we describe the main architecture decisions behind \SaiLoR.
For each decision, we state the context, the decision, and its consequences.

\subsection{Overview}\label{sec:architecture-overview}

% Overview
\SaiLoR is a desktop application built with Tauri~2~\cite{tauri2}.
The user interface is a React application in TypeScript that runs in the web view of the operating system.
It renders PDFs with pdf.js and keeps the project state in a single store.
A backend in Rust provides the native capabilities, such as file access, git, the keychain, and network requests.
The user interface reaches the backend only through commands that the backend registers.
At the time of writing, the TypeScript code base has about 50,000 lines of code without tests and about 44,000 lines of tests.
The Rust backend has about 11,500 lines in 17 modules.

\subsection{Local Files Instead of a Server}\label{sec:architecture-local}

% Local-first
Hosted tools make collaboration easy, but they keep the review data on a server (cf.~\autoref{sec:related-work-tools}).
For a research tool, this setup has drawbacks.
The data can be hard to export completely, the service can change or end, and reviewers might not be allowed to upload papers.
Therefore, we decided to keep all data in local files, following the idea of local-first software~\cite{kleppmann2019local}.
The project file and the annotation files are JSON files that people can read.
\SaiLoR keeps unknown keys in the project file, so that an older version does not delete data of a newer version.
Each annotation file names the schema version it was written under, and validation reports files with an unknown schema version.
The serialization is deterministic, and opening a project writes nothing.
As a consequence, a version control system only sees changes that a reviewer made.
The price of this decision is that \SaiLoR needs its own collaboration mechanism.

\subsection{Git as Collaboration Layer}\label{sec:architecture-git}

% Git
We decided to use git for collaboration instead of building a synchronization server.
Git is common in \SE research, offers hosting at many institutions, and keeps a full history~\cite{ram2013git}.
To keep merges small, \SaiLoR splits the annotations into one file per paper and seat.
Two reviewers who work on different seats never edit the same annotation file.

% Field-level merge
However, a line-based merge does not understand JSON.
When two people edit the same file, for example, the consolidator and a reviewer who changes the schema, git can report conflicts in lines that do not conflict in their meaning.
Therefore, \SaiLoR runs its own three-way merge on the parsed files.
It compares both sides to the common base field by field.
A side that did not change a value does not influence the result.
Repeatable entries, project settings, and schema nodes are merged separately, and answers under a removed schema field are kept.

% Safety
Using the git installation of the user brings risks.
A cloned repository can contain a configuration that runs code, for example, hooks or a file system monitor.
\SaiLoR runs git only in the backend and overrides these settings on every call, including the pager, the editor, and the \texttt{ext} protocol.
It disables interactive prompts, limits the run time of each command, and validates clone URLs, paths, and branch names.
Git operations are only allowed on repositories that the user opened or cloned in the current session.
Before \SaiLoR commits or discards selected field changes, it checks that the file on disk still matches the reviewed state, and it refuses the write otherwise.

\subsection{From Electron to Tauri}\label{sec:architecture-tauri}

% Context
Up to version 1.x, \SaiLoR was an Electron application~\cite{electron}.
The Electron main process had about 5,600 lines of TypeScript and 98 \IPC handlers.
Electron bundles its own Chromium and Node.js runtime.
As a result, the macOS installer for arm64 had a size of 93~MB, and the main process could use the full Node.js \API.

% Decision
For version 2.0.0, we moved \SaiLoR to Tauri~2~\cite{tauri2}.
Tauri uses the web view of the operating system and a backend in Rust.
In a spike, a minimal Tauri shell produced installers of about 2~MB for macOS, Windows, and the Debian and RPM packages.
The AppImage for Linux has about 81~MB, because it bundles WebKitGTK.
The shell alone is thus about 2\% of the size of the old Electron installer.
The installers of the real application are larger because they include the user interface and its assets.
% TODO: Add the final installer sizes of the 2.0.0 release per platform.
\todo{final installer sizes of 2.0.0 per platform}

% Split
Size was not the only reason.
Tauri moves the security boundary into Rust and offers a capability system and a signed updater.
We split the responsibilities as follows.
Rust owns everything that needs authorization: file input and output, process spawning, downloads, and access to the keychain.
TypeScript keeps the parsing of files, the merge logic, and the assembly of exported PDFs, because this code was already tested and does not need native access.
\SaiLoR does not expose the generic file system and shell plugins of Tauri to the user interface.
The capability file grants the web view only the default core permissions.

% Consequences
The move has consequences for users and for us.
\SaiLoR now depends on three different web views: WKWebView on macOS, WebView2 on Windows, and WebKitGTK on Linux.
We checked that pdf.js renders on all three.
The minimum macOS version is 15.
Users of version 1.x have to install version 2.0.0 manually and enter their \API keys again.
Also, not all features of the Electron version were ported yet (cf.~\autoref{sec:discussion}).

\subsection{Security Model}\label{sec:architecture-security}

% Security
A review tool opens files and repositories from other people, and it can send data to \LLM providers.
Therefore, we treat the user interface as untrusted and enforce all checks in the backend.
\SaiLoR only saves a project to a path that the user opened or chose in a native dialog during the current session.
It serves a PDF only if the file lies inside the project folder.
For a PDF outside the folder, \SaiLoR asks the user and allows only this exact path for the current session.
Exports only write to destinations that the user picked in a save dialog, and \SaiLoR refuses symbolic links as targets.
Updates are signed, and \SaiLoR downloads and installs them only when the user asks for it.

% LLM security
For \LLM providers, \SaiLoR stores \API keys in the keychain and sends a key only to the origin of its provider.
It refuses redirects for requests that carry a key.
The agent mode can fetch web pages, which opens the door to \SSRF~\cite{owasp_ssrf}.
\SaiLoR only allows \texttt{http} and \texttt{https} URLs, resolves the host name, and refuses requests to loopback, private, and other non-public addresses.
It checks every redirect again, follows at most five redirects, stops after 20 seconds, reads at most 2~MB, and only accepts text content.
The fetched text is marked as untrusted data, and the system prompt tells the model to never follow instructions from it.

\subsection{Controlled \acs{LLM} Assistance}\label{sec:architecture-llm}

% Human in the loop
Prior studies found that \LLMs can help with screening and extraction but still make errors~\cite{felizardo2024chatgpt,syriani2024screening,khraisha2024llm}.
Therefore, we designed the \LLM support so that it cannot change the review data without a human decision.
Suggestions live outside the project until the reviewer applies them.
\SaiLoR validates them against the schema and checks quoted evidence against the text of the paper.
The durable record of each apply makes the use of \LLMs visible in the data, so that authors can report it.
Treating the model as an additional reviewer seat reuses the consolidation and agreement statistics, instead of adding a separate evaluation mechanism.

% Local runtime
Some reviewers cannot send papers to a hosted provider.
For them, \SaiLoR supports local models.
For chat models, users connect a local server such as Ollama, LM Studio, llama.cpp, or vLLM.
For small decision models, \SaiLoR can manage a llama.cpp server itself~\cite{llamacpp}.
It downloads the server and the models only from an allowlist of sources, verifies their SHA-256 checksums, and starts the server on \texttt{127.0.0.1} with a random key per session.
The server uses a \GPU when one is available and the \CPU otherwise.
We verified this runtime on macOS on arm64.
The managed installation needs a llama.cpp release with a specific route for decision models.
As long as no release includes this route, the setup dialog keeps the installation disabled and points users to a server that they run themselves.
